\chapter{Регистры и счётчики}
\label{ch:regs}

Один триггер хранит один бит. Соединив несколько триггеров, получают узлы,
которые хранят числа, сдвигают их и считают.

\newcommand{\dff}[3]{% имя, x, y  -- D-триггер 12x14 мм
  \draw[thick, fill=green!8] (#2,#3) rectangle ++(1.2,1.4);
  \node[font=\sffamily\scriptsize] at (#2+0.22,#3+1.05) {D};
  \node[font=\sffamily\scriptsize] at (#2+0.98,#3+1.05) {Q};
  \draw[thick] (#2,#3+0.2) -- (#2+0.18,#3+0.35) -- (#2,#3+0.5);
  \coordinate (#1-d) at (#2,#3+1.05); \coordinate (#1-q) at (#2+1.2,#3+1.05);
  \coordinate (#1-c) at (#2,#3+0.35);}

\section{Регистр хранения}

\term{Регистр} --- это группа D-триггеров с общим тактовым сигналом. По
фронту все триггеры одновременно запоминают свои входы, и регистр хранит
$n$-разрядное число до следующего фронта. На схемах регистр обозначают
буквами RG.

\begin{figure}[H]
\centering
\begin{tikzpicture}
  \foreach \i in {0,...,3} {
    \pgfmathsetmacro{\x}{(3-\i)*2.3}
    \dff{r\i}{\x}{0}
    \draw[wire] (r\i-d) -- ++(-0.4,0) |- ($(r\i-d)+(-0.4,1.0)$) node[above] {$d_\i$};
    \draw[wire] (r\i-q) -- ++(0.35,0) -- ++(0,-1.5) node[below] {$q_\i$};
    \draw[wire, FpgaBlue] (r\i-c) -- ++(-0.3,0) -- ++(0,-0.75);
  }
  \draw[wire, FpgaBlue] (-1.4,-0.4) node[left, text=FpgaDark] {CLK} -- (6.6,-0.4);
  \foreach \i in {0,...,3} { \pgfmathsetmacro{\x}{(3-\i)*2.3-0.3} \fill[FpgaBlue] (\x,-0.4) circle (1.3pt); }
\end{tikzpicture}
\caption{Четырёхразрядный регистр: четыре D-триггера с общим тактом}
\label{fig:reg}
\end{figure}

\begin{figure}[H]
\centering
\begin{tikztimingtable}[timing/slope=0.08, xscale=1.5, yscale=1.4, timing/font=\sffamily\small, timing/rowdist=3.2ex]
  CLK      & 8{1L 1H} \\
  $D[3:0]$ & 2D{5} 1.5D{9} 1D{7} 1.5D{3} 4D{C} 6D{0} \\
  $Q[3:0]$ & 1D{?} 2D{5} 2D{9} 2D{3} 4D{C} 5D{0} \\
\end{tikztimingtable}
\caption{Регистр «фотографирует» значение шины $D$ в моменты передних фронтов
CLK. Значение 7 продержалось на входе меньше такта и пришлось между
фронтами --- регистр его просто не заметил}
\end{figure}

Регистры --- это «переменные» аппаратуры. В процессоре в регистрах хранятся
операнды и результаты вычислений; в любом устройстве на ПЛИС --- текущее
состояние, счётчики, данные между этапами обработки. Чтобы регистр менял
значение не на каждом фронте, а только по команде, используют триггеры с
входом разрешения CE (рис.~\ref{fig:ce}).

\section{Сдвиговый регистр}

Соединим триггеры цепочкой: выход каждого подаётся на вход следующего. По
каждому фронту содержимое \term{сдвигается} на один разряд, а в первый
триггер записывается новый бит с входа $D_\text{in}$.

\begin{figure}[H]
\centering
\begin{tikzpicture}
  \foreach \i in {0,...,3} {
    \pgfmathsetmacro{\x}{\i*2.3}
    \dff{s\i}{\x}{0}
    \draw[wire, FpgaBlue] (s\i-c) -- ++(-0.3,0) -- ++(0,-0.75);
  }
  \draw[wire] (-1.0,1.05) node[left] {$D_\text{in}$} -- (s0-d);
  \foreach \i [evaluate=\i as \j using int(\i+1)] in {0,1,2} {
    \draw[wire] (s\i-q) -- (s\j-d);
    \fill ($(s\i-q)+(0.3,0)$) circle (1.3pt);
    \draw[wire] ($(s\i-q)+(0.3,0)$) -- ++(0,1.0) node[above] {$q_\i$};
  }
  \draw[wire] (s3-q) -- ++(0.7,0) node[right] {$q_3 = D_\text{out}$};
  \draw[wire, FpgaBlue] (-1.0,-0.4) node[left, text=FpgaDark] {CLK} -- (6.6,-0.4);
\end{tikzpicture}
\caption{Четырёхразрядный сдвиговый регистр}
\label{fig:shift}
\end{figure}

\begin{figure}[H]
\centering
\begin{tikztimingtable}[timing/slope=0.08, xscale=1.5, yscale=1.4, timing/font=\sffamily\small, timing/rowdist=3.2ex]
  CLK            & 8{1L 1H} \\
  $D_\text{in}$  & 2H 2L 12L \\
  $q_0$          & 1L 2H 13L \\
  $q_1$          & 3L 2H 11L \\
  $q_2$          & 5L 2H 9L \\
  $q_3$          & 7L 2H 7L \\
\end{tikztimingtable}
\caption{Единица, поданная на вход, продвигается по регистру на один разряд
за такт}
\end{figure}

\subsection{Зачем сдвигать}

\begin{description}
  \item[Последовательный $\to$ параллельный код.] Биты приходят по одному
    проводу друг за другом (так работают интерфейсы UART, SPI, I$^2$C). После
    $n$ тактов все $n$ битов оказываются в регистре одновременно и могут быть
    прочитаны как число.
  \item[Параллельный $\to$ последовательный код.] Обратная задача при передаче:
    число загружается в регистр целиком, а затем «выталкивается» по одному
    биту за такт. Для загрузки перед каждым триггером ставят мультиплексор
    «сдвиг/загрузка».
  \item[Задержка сигнала] на заданное число тактов.
  \item[Умножение и деление на 2]: сдвиг двоичного числа влево на один разряд
    умножает его на 2, вправо --- делит.
\end{description}

\begin{figure}[H]
\centering
\begin{tikzpicture}[bit/.style={draw=FpgaBlue, fill=FpgaLight, minimum size=8mm, font=\ttfamily}]
  \node[font=\sffamily\small, anchor=east] at (-0.6,0) {передатчик};
  \foreach \b [count=\i from 0] in {1,0,1,1} \node[bit] at (\i*0.8,0) {\b};
  \draw[very thick, FpgaOrange, -{Stealth}] (3.0,0) -- node[above, font=\sffamily\small, text=FpgaDark] {1 провод: \code{1}, \code{1}, \code{0}, \code{1}\ldots} (7.4,0);
  \foreach \b [count=\i from 0] in {1,0,1,1} \node[bit] at (7.8+\i*0.8,0) {\b};
  \node[font=\sffamily\small, anchor=west] at (10.8,0) {приёмник};
  \node[lbl] at (1.2,-0.8) {параллельный $\to$ послед.};
  \node[lbl] at (9.0,-0.8) {послед. $\to$ параллельный};
\end{tikzpicture}
\caption{Передача числа по одному проводу с помощью двух сдвиговых регистров}
\end{figure}

\subsection{Кольцевой регистр и генератор псевдослучайных чисел}

Если выход последнего триггера соединить со входом первого, получится
\term{кольцевой регистр}: записанная в него единица будет бесконечно бегать по
кругу (эффект «бегущего огня» в главе~\ref{ch:examples}).

А если на вход подать «Исключающее ИЛИ» от нескольких разрядов, получится
\term{регистр сдвига с линейной обратной связью} (LFSR). При правильном выборе
отводов он перебирает все $2^n-1$ ненулевых состояний в «перемешанном»
порядке. Такие регистры используются как простые генераторы
псевдослучайных чисел, в помехоустойчивом кодировании и для шифрования.

\begin{figure}[H]
\centering
\begin{minipage}{0.55\textwidth}
\centering
\begin{tikzpicture}
  \foreach \i in {0,...,2} {
    \pgfmathsetmacro{\x}{\i*2.0}
    \dff{l\i}{\x}{0}
  }
  \draw[wire] (l0-q) -- (l1-d); \draw[wire] (l1-q) -- (l2-d);
  \node[gXOR, minimum width=9mm, minimum height=7mm, rotate=180] (x) at (2.6,2.3) {};
  \draw[wire] (l2-q) -- ++(0.4,0) |- (x.input 1);
  \fill ($(l1-q)+(0.3,0)$) circle (1.3pt);
  \draw[wire] ($(l1-q)+(0.3,0)$) |- (x.input 2);
  \draw[wire] (x.output) -| ($(l0-d)+(-0.4,0)$) -- (l0-d);
  \foreach \i in {0,...,2} \node[lbl, below=1mm] at ($(l\i-c)+(0.6,-0.4)$) {$q_\i$};
\end{tikzpicture}
\end{minipage}\hfill
\begin{minipage}{0.4\textwidth}
\centering
\small
\begin{tabular}{c|ccc}
\toprule
такт & $q_0$ & $q_1$ & $q_2$ \\
\midrule
0 & 1 & 0 & 0 \\
1 & 0 & 1 & 0 \\
2 & 1 & 0 & 1 \\
3 & 1 & 1 & 0 \\
4 & 1 & 1 & 1 \\
5 & 0 & 1 & 1 \\
6 & 0 & 0 & 1 \\
7 & 1 & 0 & 0 \\
\bottomrule
\end{tabular}
\end{minipage}
\caption{Трёхразрядный LFSR ($q_0 \leftarrow q_1\oplus q_2$) перебирает все семь
ненулевых состояний и повторяется}
\end{figure}

\section{Счётчики}

\term{Счётчик} --- регистр, значение которого по каждому такту
увеличивается (или уменьшается) на единицу. Счётчики измеряют время, делят
частоту, перебирают адреса памяти.

\subsection{Асинхронный счётчик}

Посмотрите, как меняются разряды при счёте $0, 1, 2, \ldots, 7$: младший бит
переключается каждый шаг, следующий --- в два раза реже, следующий --- ещё в
два раза реже. Каждый разряд переключается, когда предыдущий переходит из 1
в 0. Значит, можно соединить T-триггеры цепочкой: тактовый вход каждого
следующего триггера подключить к \emph{инверсному} выходу предыдущего.

\begin{figure}[H]
\centering
\begin{tikzpicture}
  \foreach \i in {0,...,2} {
    \pgfmathsetmacro{\x}{\i*2.8}
    \draw[thick, fill=green!8] (\x,0) rectangle ++(1.2,1.4);
    \node[font=\sffamily\scriptsize] at (\x+0.22,1.05) {T};
    \node[font=\sffamily\scriptsize] at (\x+0.98,1.05) {Q};
    \node[font=\sffamily\scriptsize] at (\x+0.98,0.35) {$\overline{\text{Q}}$};
    \draw[thick] (\x,0.2) -- (\x+0.18,0.35) -- (\x,0.5);
    \draw[wire] (\x,1.05) -- ++(-0.3,0) node[left, font=\scriptsize] {1};
    \draw[wire] (\x+1.2,1.05) -- ++(0.4,0) -- ++(0,0.8) node[above] {$q_\i$};
  }
  \draw[wire, FpgaBlue] (-1.0,0.35) node[left, text=FpgaDark] {CLK} -- (0,0.35);
  \draw[wire, FpgaRed] (1.2,0.35) -- (2.8,0.35);
  \draw[wire, FpgaRed] (4.0,0.35) -- (5.6,0.35);
\end{tikzpicture}
\caption{Асинхронный (последовательный) трёхразрядный счётчик на T-триггерах}
\label{fig:ripple}
\end{figure}

\begin{figure}[H]
\centering
\begin{tikztimingtable}[timing/slope=0.05, xscale=1.2, yscale=1.4, timing/font=\sffamily\small, timing/rowdist=3.2ex]
  CLK     & 16{1L 1H} \\
  $q_0$   & 1L 7{2H 2L} 2H 1L \\
  $q_1$   & 3L 3{4H 4L} 4H 1L \\
  $q_2$   & 7L 8H 8L 8H 1L \\
  число   & 1D{0} 2D{1} 2D{2} 2D{3} 2D{4} 2D{5} 2D{6} 2D{7} 2D{0} 2D{1} 2D{2} 2D{3} 2D{4} 2D{5} 2D{6} 2D{7} 1D{0} \\
\end{tikztimingtable}
\caption{Каждый следующий разряд меняется вдвое реже: счётчик перебирает числа
$0\ldots7$, а разряд $q_k$ имеет частоту $f_\text{CLK}/2^{k+1}$}
\label{fig:count_t}
\end{figure}

Недостаток асинхронного счётчика --- триггеры переключаются не одновременно,
а «волной», каждый с задержкой относительно предыдущего. В момент перехода
$011 \to 100$ на выходе на короткое время появляются ложные значения
($011\to010\to000\to100$). В длинном счётчике эта «волна» идёт долго.

\subsection{Синхронный счётчик}

В синхронном счётчике все триггеры тактируются \emph{одним} сигналом, а новое
значение вычисляет комбинационная схема --- сумматор, прибавляющий единицу.
Все разряды меняются одновременно по фронту. Именно так строятся все
счётчики в ПЛИС.

\begin{figure}[H]
\centering
\begin{tikzpicture}
  \node[draw=FpgaGreen, fill=green!8, thick, minimum width=20mm, minimum height=18mm, font=\sffamily\small, align=center] (reg) at (4.4,0) {регистр\\\scriptsize $n$ бит};
  \node[font=\sffamily\scriptsize, anchor=west] at (reg.west) {D};
  \node[font=\sffamily\scriptsize, anchor=east] at (reg.east) {Q};
  \draw[thick] ([xshift=-1.2mm]reg.south) -- ([yshift=1.8mm]reg.south) -- ([xshift=1.2mm]reg.south);
  \draw[wire, FpgaBlue] (reg.south) -- ++(0,-0.8) node[below, text=FpgaDark] {CLK};
  \node[draw=FpgaBlue, fill=FpgaLight, thick, circle, minimum size=12mm, font=\sffamily\bfseries] (add) at (1.4,0) {$+1$};
  \draw[bus, -{Stealth}] (add.east) -- (reg.west);
  \draw[bus, -{Stealth}] (reg.east) -- ++(1.8,0) node[right, text=FpgaDark] {$Q$};
  \draw[bus, -{Stealth}] ($(reg.east)+(0.9,0)$) -- ++(0,1.6) -| (add.north);
  \fill[FpgaBlue] ($(reg.east)+(0.9,0)$) circle (1.8pt);
  \node[lbl] at (4.4,1.9) {обратная связь};
\end{tikzpicture}
\caption{Синхронный счётчик: регистр + сумматор. Каждый такт регистр
записывает $Q+1$}
\label{fig:sync_counter}
\end{figure}

\subsection{Делитель частоты и счётчик по модулю}

Как видно на рис.~\ref{fig:count_t}, каждый разряд счётчика --- это
делитель частоты на степень двойки. А как поделить на число, не являющееся
степенью двойки, например на 10? Нужен \term{счётчик по модулю $N$}: он считает
$0, 1, \ldots, N-1$ и затем сбрасывается в 0. Для этого к схеме добавляют
компаратор, который сравнивает значение с $N-1$, и мультиплексор, который
вместо $Q+1$ подаёт на регистр 0.

\begin{figure}[H]
\centering
\begin{tikzpicture}
  \node[draw=FpgaGreen, fill=green!8, thick, minimum width=18mm, minimum height=16mm, font=\sffamily\small] (reg) at (5.6,0) {регистр};
  \draw[wire, FpgaBlue] (reg.south) -- ++(0,-0.6) node[right, text=FpgaDark, font=\sffamily\small] {CLK};
  \draw[thick, fill=orange!15] (3.0,0.9) -- (3.8,0.5) -- (3.8,-0.5) -- (3.0,-0.9) -- cycle;
  \node[font=\sffamily\scriptsize] at (3.4,0) {MUX};
  \node[font=\tiny, anchor=west] at (3.02,0.5) {0}; \node[font=\tiny, anchor=west] at (3.02,-0.5) {1};
  \node[draw=FpgaBlue, fill=FpgaLight, thick, circle, minimum size=10mm, font=\sffamily\bfseries\small] (add) at (1.2,0.5) {$+1$};
  \draw[bus, -{Stealth}] (add.east) -- (3.0,0.5);
  \draw[bus, -{Stealth}] (1.8,-0.5) node[left, text=FpgaDark] {0} -- (3.0,-0.5);
  \draw[bus, -{Stealth}] (3.8,0) -- (reg.west);
  \draw[bus, -{Stealth}] (reg.east) -- ++(1.6,0) node[right, text=FpgaDark] {$Q$};
  \draw[bus] ($(reg.east)+(0.8,0)$) -- ++(0,1.6) -| (add.north);
  \fill[FpgaBlue] ($(reg.east)+(0.8,0)$) circle (1.8pt);
  \node[block, minimum width=20mm] (cmp) at (5.6,-2.9) {$Q = N-1$?};
  \draw[bus, -{Stealth}] ($(reg.east)+(0.8,0)$) |- (cmp.east);
  \draw[arr] (cmp.west) -| node[pos=0.3, below, lbl] {сброс} (3.4,-0.72);
\end{tikzpicture}
\caption{Счётчик по модулю $N$}
\end{figure}

\begin{figure}[H]
\centering
\begin{tikzpicture}
\begin{axis}[
  width=0.82\textwidth, height=5.6cm,
  xlabel={Разрядность счётчика $n$}, ylabel={Частота, Гц},
  ymode=log, ymin=0.1, ymax=3e7, xmin=0, xmax=26, xtick={1,4,8,12,16,20,25},
  label style={font=\sffamily\small}, tick label style={font=\sffamily\small},
  axis lines*=left, grid=major, grid style={FpgaGray!20},
]
  \addplot[very thick, FpgaBlue, mark=*, mark size=1.5pt, samples at={1,2,...,25}] {27e6/2^x};
  \addplot[FpgaRed, thick, dashed, domain=0:26, samples=2] {1};
  \node[font=\sffamily\small, text=FpgaRed, anchor=south west] at (axis cs:1,1.3) {1 Гц --- мигание, заметное глазу};
  \node[font=\sffamily\small, text=FpgaBlue, anchor=south west] at (axis cs:25,0.8) {$n=25$};
\end{axis}
\end{tikzpicture}
\caption{Частота на выходе старшего разряда $n$-битного счётчика при такте
27 МГц: $f = 27\,\text{МГц}/2^n$. Чтобы получить мигание около 1 Гц, нужен
счётчик примерно на 25 бит}
\label{fig:div_plot}
\end{figure}

\begin{ideabox}
\textbf{Регистр} --- группа триггеров с общим тактом, хранит число.
\textbf{Сдвиговый регистр} перемещает биты по цепочке и преобразует
последовательный код в параллельный и обратно. \textbf{Синхронный счётчик}
--- регистр с сумматором в обратной связи; его разряды делят частоту такта
на степени двойки.
\end{ideabox}
